
\chapter{黎曼 \texorpdfstring{$\zeta$}{Zeta} 函数初探}\label{ch:riemann_zeta_intro}

\vspace{0.6em}
\begin{quote}
\small
\[
  \sum_{n=1}^{\infty}\frac{1}{n^{s}}
  \;=\;
  \prod_{p}\frac{1}{1-p^{-s}}
  \qquad (s>1).
\]
{\raggedright
素数经由无穷级数与无穷乘积进入分析。\\[2pt]
\hfill——L.~Euler，\textit{Variae observationes circa series infinitas}（1737）}
\end{quote}
\vspace{0.6em}

\section{引言}\label{sec:intro_zeta_intro}

$\zeta(s)$ 的 Dirichlet 级数表达式是全体正整数上的函数，而 $\zeta(s)$ 的 Euler 乘积的数学表达式是全体素数上的函数；
二者在形式上就已经蕴含了与素数分布及其他素数问题的天然内在联系，使 $\zeta$ 成为解析数论最重要的函数。
Euler 于 1734--1735 年解决著名的巴塞尔问题，得到实轴上
$\zeta(2)=\pi^2/6$；又于 1737 年在
\textit{Variae observationes circa series infinitas} 中系统导出 Euler 乘积
$\sum n^{-s}=\prod_p(1-p^{-s})^{-1}$（$s>1$）。
黎曼（1859）将自变量推广到复数，导出了素数分布的黎曼显式公式，并指出非平凡零点控制素数分布及其渐近误差。

本章只处理半平面 $\operatorname{Re}(s)>1$ 上的原初定义。
逻辑主线为：Dirichlet 级数的收敛、Euler 乘积（与级数恒等）、
两种表达式的分工、$\operatorname{Re}s>1$ 无零点、全纯性；
\textbf{其后}由 Euler 乘积建立 $\mu$、$\Lambda$ 的 Dirichlet 生成函数
（$1/\zeta$、$-\zeta'/\zeta$、$\log\zeta$）；
第~\ref{ch:discrete_continuous}~章再把这些生成函数代入第~\ref{ch:integral_transforms}~章的 Perron 积分公式，
把部分和 $M,\psi,J$ 写成竖线积分。
再述共轭对称，并以截断级数示意与局限说明级数边界。
完整解析延拓见第~\ref{ch:riemann_zeta_continuation}~章。

%==============================================================
\section{\texorpdfstring{$\zeta$}{Zeta(s)} 函数的定义}
%==============================================================
\subsection{Dirichlet 级数的定义}

\begin{definition}{一般 Dirichlet 级数}{}
形如
\begin{equation}
    F(s) =\sum_{n=1}^{\infty} \frac{a_n}{n^s},
    \qquad s = \sigma + it \in \mathbb{C}
\end{equation}
的级数称为\textbf{（一般）Dirichlet 级数}，其中
$\{a_n\}_{n=1}^{\infty}$ 为复数列，称为该级数的\textbf{系数序列}。
\end{definition}

特别地，当 $ a_n = 1 $ 对所有 $ n $ 成立时，该级数即为黎曼 $\zeta(s)$的 Dirichlet 级数形式：
\begin{equation}
  \zeta(s) =\sum_{n=1}^{\infty} n^{-s} =\sum_{n=1}^{\infty} \frac{1}{n^s}.
\end{equation}

\begin{theorem}{绝对收敛半平面}{convergence_halfplane}
$\zeta(s)$ 的 Dirichlet 级数 $\sum_{n=1}^{\infty} n^{-s}$ 在半平面 $\operatorname{Re}(s) > 1$ 上\textbf{绝对收敛}。
具体而言，当 $\sigma = \operatorname{Re}(s) > 1$ 时，
\begin{equation}
\sum_{n=1}^{\infty} \bigl| n^{-s} \bigr|
  =\sum_{n=1}^{\infty} n^{-\sigma}
  < \infty.
\end{equation}
因而其\textbf{绝对收敛横坐标}为
\[
  \sigma_a = 1.
\]
\end{theorem}

\begin{proof}
因为
\[
  \bigl| n^{-s} \bigr|
  = \bigl| e^{-s \log n} \bigr|
  = e^{-\sigma \log n}
  = n^{-\sigma},
\]
于是 $\sum |n^{-s}|=\sum n^{-\sigma}$，这是参数为 $\sigma$ 的 p-级数（即 $\sum 1/n^{\sigma}$）。
当 $\sigma > 1$ 时该级数收敛，故原级数绝对收敛；
当 $\sigma \le 1$ 时 $\sum n^{-\sigma}$ 发散，故绝对收敛不能越过 $\operatorname{Re}s=1$，即 $\sigma_a=1$。
\end{proof}

又，在 $s=1$ 处级数化为调和级数 $\sum 1/n$ 而发散，故收敛横坐标满足 $\sigma_c\ge 1$；
结合 $\sigma_c\le\sigma_a=1$，得\textbf{收敛横坐标} $\sigma_c=1$。
本书在 $\operatorname{Re}s>1$ 上用绝对收敛的 Dirichlet 级数定义 $\zeta(s)$
（全纯性见第~\ref{sec:holo}~节）；$\operatorname{Re}s\le 1$ 留给后文解析延拓。
绝对收敛保证可重排，从而可导出 Euler 乘积（依赖唯一分解）。

%==============================================================
\section{Euler 乘积及其适用范围}\label{sec:euler-product}
%==============================================================

上一节用 Dirichlet 级数在 $\operatorname{Re}s>1$ 上定义了 $\zeta(s)$。
本节证明：在同一半平面上，该级数与 Euler 乘积
$\prod_p(1-p^{-s})^{-1}$ \textbf{恒等}
（绝对收敛下的重排与算术基本定理；筛法为直观版）。
后半说明乘积在 $\sigma=1$ 上的失效。
两种表达式既已建立，其总结与后文分工见下一节。

\subsection{\texorpdfstring{$\zeta(s)$}{Zeta(s)} Euler 乘积的推导}

Euler 乘积给出 $\zeta(s)$ 与素数之间的直接联系。
Euler 于 1737 年的 \textit{Variae observationes circa series infinitas} 中系统写出该乘积；
此前（1734--1735）他已在解决巴塞尔问题时得到实轴特殊值 $\zeta(2)=\pi^2/6$。

\begin{theorem}{Euler 乘积}{euler_product}
对 $\operatorname{Re}(s) > 1$，有
\begin{equation}
  \boxed{
    \zeta(s) =\prod_{p \text{ 素数}} \frac{1}{1 - p^{-s}}
             =\prod_{p} \left(1 - \frac{1}{p^s}\right)^{-1}.
  }
\end{equation}
乘积遍历所有素数 $p = 2, 3, 5, 7, 11, \ldots$
\end{theorem}

\begin{proof}[基于唯一素因子分解的证明]
当 $\operatorname{Re}(s) > 1$ 时，Dirichlet 级数 $\sum n^{-s}$ 绝对收敛，可按素数分解重排求和项。
由算术基本定理，每个正整数 $n \geq 1$ 可唯一写成素数幂的乘积 $n =\prod_p p^{k_p}$（其中 $k_p \geq 0$ 且仅有限个 $k_p > 0$）。
因此
\[
\zeta(s) =\sum_{n=1}^\infty n^{-s} =\sum_{n=1}^\infty\prod_p p^{-k_p s}.
\]
绝对收敛允许交换求和顺序，故上式等于
\[\prod_p \left(\sum_{k=0}^\infty p^{-k s} \right).
\]
每个内层级数为几何级数（因为 $|p^{-s}| = p^{-\sigma} < 1$），求和得 $(1 - p^{-s})^{-1}$。
于是
\[
\zeta(s) =\prod_p (1 - p^{-s})^{-1}.
\]
\end{proof}


同一公式亦可由 Euler 式的逐步筛除直观推出（类比埃拉托斯特尼筛法），有助于理解素数与级数的联系。

考虑级数 $\zeta(s) =\sum_{n=1}^\infty n^{-s}$。首先乘以 $(1 - 2^{-s})$：
\[
(1 - 2^{-s})\zeta(s) =\sum_{n=1}^\infty n^{-s} -\sum_{n=1}^\infty (2n)^{-s} =\sum_{\substack{n=1 \\ 2 \nmid n}}^\infty n^{-s},
\]
即筛去了所有含素数因子 $2$ 的项（偶数项）。

接着乘以 $(1 - 3^{-s})$：
\[
(1 - 3^{-s})(1 - 2^{-s})\zeta(s) =\sum_{\substack{n=1 \\ 2,3 \nmid n}}^\infty n^{-s},
\]
筛去了含 $3$ 的项。依此类推，对每个素数 $p$ 乘以 $(1 - p^{-s})$，逐步筛去所有含该素数因子的项。

当我们依次对所有素数进行这一操作（理论上需进行无穷步），剩余的求和项仅为 $n=1$（其系数为 $1^{-s} = 1$），因为任何 $n > 1$ 都至少含有一个素数因子，最终被筛除。因此
\[
\left(\prod_p (1 - p^{-s}) \right) \zeta(s) = 1,
\]
从而得到
\[
\zeta(s) =\prod_p (1 - p^{-s})^{-1}.
\]
在 $\operatorname{Re}(s) > 1$ 时，上述过程可通过有限素数乘积的极限严格化（利用绝对收敛性控制余项），与前述证明等价。

\subsection{Euler 乘积的适用范围与 \texorpdfstring{$\sigma=1$}{sigma=1} 处的失效}
\label{sec:euler-fail}
\label{sec:euler-limitations}

本节只讨论\textbf{乘积工具}本身：它在何处与级数等价，又在何处不能再被当作证明手段。
（以级数定义 $\zeta$ 的全局不足——临界带、特殊值、函数方程、数值伪像——见后文第~\ref{sec:limitations}~节。）

\paragraph{与级数的相等限于 $\operatorname{Re}s>1$}
定理~\ref{thm:euler_product} 依赖绝对收敛以交换求和顺序。
因而
\[
  \sum_{n=1}^{\infty}n^{-s}
  =\prod_{p}(1-p^{-s})^{-1}
\]
是 $\operatorname{Re}s>1$ 上的恒等式，不是全平面上的定义。

\paragraph{在直线 $\sigma=1$ 上绝对收敛失效}
一旦落到 $s=1+it$（$t$ 任意实数），乘积
$\prod_p (1-p^{-1-it})^{-1}$ \textbf{不再绝对收敛}：
当 $p\to\infty$ 时
\begin{equation}
\left|\frac{p^{-1-it}}{1-p^{-1-it}}\right| = \frac{1}{p|1-p^{-1-it}|} \sim \frac{1}{p},
\end{equation}
而由下面的 Mertens 定理，$\sum_p 1/p=+\infty$，故
$\sum_p \bigl|\frac{p^{-1-it}}{1-p^{-1-it}}\bigr|=+\infty$。

\begin{theorem}{Mertens 定理（素数倒数和）}{mertens_harmonic}
素数的倒数和发散，且有渐近
\begin{equation}
  \sum_{p\le x}\frac{1}{p}
  \;=\;
  \log\log x + C + o(1)
  \qquad(x\to\infty),
  \label{eq:mertens_sum_1_p}
\end{equation}
其中 $C$ 为 Mertens 常数。特别地 $\sum_p 1/p=+\infty$。
\end{theorem}

\begin{proof}[纲要（仅证 $\sum_p 1/p=+\infty$）]
对实数 $\sigma>1$，由 Euler 乘积取对数，
\[
  \log\zeta(\sigma)
  = -\sum_p\log(1-p^{-\sigma})
  = \sum_p\sum_{k\ge 1}\frac{p^{-k\sigma}}{k}
  = \sum_p p^{-\sigma} + R(\sigma),
\]
其中 $0\le R(\sigma)\le\sum_p\sum_{k\ge 2}p^{-k\sigma}\le\sum_p\frac{p^{-2\sigma}}{1-p^{-\sigma}}$ 在 $\sigma\to 1^+$ 时保持有界。
又对任意固定 $N$，当 $\sigma\to 1^+$ 时 $\zeta(\sigma)\ge\sum_{n\le N}n^{-\sigma}\to H_N$，
故 $\zeta(\sigma)\to+\infty$，从而 $\log\zeta(\sigma)\to+\infty$。
于是 $\sum_p p^{-\sigma}\to+\infty$（$\sigma\to 1^+$）。
由单调收敛（$p^{-\sigma}\uparrow p^{-1}$）即得 $\sum_p 1/p=+\infty$。

更精细的 $\sum_{p\le x}1/p=\log\log x+C+o(1)$ 需分部求和等，
见 Apostol 或 Titchmarsh；本书此处只用发散性。
\end{proof}


因此后文无零点定理\textbf{仅适用于 $\operatorname{Re}s>1$}：
$\sigma=1$、$t\neq 0$ 时 $\zeta(1+it)\neq 0$ 不能再用 Euler 乘积直接推出，
须在解析延拓之后另证（Hadamard 三角不等式等，见第~\ref{ch:zero_distribution}~章）。

\paragraph{对本节的定位}
在 $\operatorname{Re}s>1$ 内，Euler 乘积是连接素数与 $\zeta$ 的最清晰工具。

%==============================================================
\section{级数形式与乘积形式：恒等与分工}
\label{sec:series_product_roles}
%==============================================================

\subsection{同一函数的两种表达式}

在 $\operatorname{Re}s>1$ 上，我们已有：
\begin{enumerate}[label=(\roman*)]
  \item \textbf{级数形式}（定义）
        \[
          \zeta(s)=\sum_{n=1}^{\infty}n^{-s};
        \]
  \item \textbf{乘积形式}（第~\ref{sec:euler-product}~节定理~\ref{thm:euler_product}）
        \[
          \zeta(s)=\prod_{p}\bigl(1-p^{-s}\bigr)^{-1}.
        \]
\end{enumerate}
二者\textbf{恒等}：左边由绝对收敛的 Dirichlet 级数定义 $\zeta$，右边由 Euler 乘积给出；
证明表明乘积不是另一个函数，而是级数在唯一分解下的重写。
因而在 $\operatorname{Re}s>1$ 上谈论「级数的 $\zeta$」与「乘积的 $\zeta$」是同一对象。

恒等\textbf{仅}在该半平面上成立。出了 $\operatorname{Re}s>1$，级数发散，不能再把上述乘积当作原定义的自动延拓
（乘积在 $\sigma=1$ 上的绝对收敛失效见第~\ref{sec:euler-fail}~节；
以级数为定义的全局不足见第~\ref{sec:limitations}~节）。

\subsection{后文如何分别使用两种形式}

恒等之后，可按问题选用更方便的写法，并在需要时用恒等式翻译到另一写法。

\paragraph{乘积形式的用武之地}
乘积的因子明确跑遍素数，适合处理与素数结构直接相关的问题：
\begin{itemize}
  \item \textbf{无零点}（下一节）：每个因子 $(1-p^{-s})^{-1}\neq 0$，
        且绝对收敛无穷乘积在无零因子时整体非零，故 $\operatorname{Re}s>1$ 时 $\zeta(s)\neq 0$；
  \item \textbf{生成函数}（再后）：对乘积取倒数得 $1/\zeta=\sum\mu n^{-s}$，
        取对数再求导得 $-\zeta'/\zeta=\sum\Lambda n^{-s}$，
        取对数得 $\log\zeta=\sum(\Lambda/\log n)n^{-s}$，
        从而把第~\ref{ch:elementary_number_theoretic_functions}~章的 $\mu$、$\Lambda$ 接到 $\zeta$ 上。
\end{itemize}

\paragraph{级数形式的用武之地}
级数是通项 $n^{-s}$ 的叠加，适合作分析分解与数值截断：
\begin{itemize}
  \item \textbf{实部与虚部}、C--R 方程与全纯性：
        将 $n^{-s}=n^{-\sigma}e^{-it\log n}$ 拆成余弦/正弦级数；
  \item \textbf{共轭对称} $\zeta(\bar s)=\overline{\zeta(s)}$：对级数逐项取共轭（或 Schwarz 反射）；
  \item \textbf{截断示意}：有限和 $\zeta_N=\sum_{n\le N}n^{-s}$ 便于作实部、虚部、模曲面。
\end{itemize}

\paragraph{阅读指引}
\begin{center}
\renewcommand{\arraystretch}{1.35}
\begin{tabular}{|l|l|l|}
\hline
\textbf{内容} & \textbf{主要使用} & \textbf{位置} \\
\hline
无零点（$\operatorname{Re}s>1$） & 乘积 & 下一节 \\
实虚分解、全纯 & 级数 & 再后 \\
共轭对称、几何示意 & 级数 & 再后 \\
$\mu$、$\Lambda$ 生成函数 & 乘积 & 再后 \\
\hline
\end{tabular}
\end{center}

\vspace{0.4em}
\noindent
以下先用乘积证明无零点，再用级数完成分析性质，最后回到乘积建立生成函数。

%==============================================================
\section{\texorpdfstring{$\zeta(s)$}{Zeta(s)} 在 \texorpdfstring{$\operatorname{Re}(s) > 1$}{Re(s) > 1} 区域的无零点性}
%==============================================================

由 Euler 乘积可直接推出，$\zeta(s)$ 在 $\operatorname{Re}(s) > 1$ 内没有零点。
这一性质为后文模曲面在 $\sigma > 1$ 上「无触底」的观察提供理论基础。

\begin{lemma}{绝对收敛无穷乘积的非零性}{}
\label{lem:product}
设 $\{a_n\}$ 为复数列，若 $\sum_{n=1}^\infty |a_n| < \infty$，
则无穷乘积 $\prod_{n=1}^\infty (1+a_n)$ 绝对收敛，
且为零当且仅当某个因子 $1+a_n = 0$。
\end{lemma}

\begin{theorem}{无零点定理}{}
当 $\operatorname{Re}(s) > 1$ 时，
\[
\zeta(s) \neq 0.
\]
\end{theorem}

\begin{proof}
根据 Euler 乘积，当 $\operatorname{Re}(s) > 1$ 时，
\[
\zeta(s) =\prod_{p \text{ 素数}} (1 - p^{-s})^{-1}.
\]
对任意素数 $p$ 和 $\sigma = \operatorname{Re}(s) > 1$，我们有
\[
|p^{-s}| = p^{-\sigma} < 1,
\]
因此 $1 - p^{-s} \neq 0$，进而每个 Euler 因子 $(1 - p^{-s})^{-1}$ 均非零。

为验证无穷乘积非零，记 $a_p = p^{-s}/(1-p^{-s})$，则 $(1-p^{-s})^{-1}=1+a_p$。由于 $\sigma > 1$ 时对所有素数 $p$ 均有 $p^{-\sigma} \leq 2^{-\sigma} < \frac{1}{2}$，从而 $|1-p^{-s}| \geq 1 - |p^{-s}| = 1 - p^{-\sigma} > \frac{1}{2}$，因此有：
\begin{equation}
|a_p| = \frac{|p^{-s}|}{|1-p^{-s}|} = \frac{p^{-\sigma}}{|1-p^{-s}|} \leq 2 p^{-\sigma},
\end{equation}
进而
\begin{equation}\sum_p |a_p|
    \leq 2 \sum_p p^{-\sigma}
    \leq 2 \sum_{n=2}^\infty n^{-\sigma}
    = 2(\zeta(\sigma) - 1) < \infty,
    \qquad \sigma > 1,
\end{equation}
从而引理 \ref{lem:product} 的条件满足，无穷乘积绝对收敛且非零，故
\[
\zeta(s) \neq 0.
\]
\end{proof}


注意：本节论证完全依赖 Euler 乘积在 $\sigma > 1$ 的绝对收敛，故\textbf{无零点定理仅适用于 $\sigma > 1$}。Euler 乘积在 $\sigma = 1$ 处不绝对收敛（见第 \ref{sec:euler-fail} 节），因而该证明不能直接推广到 $\sigma = 1$。边界线 $\zeta(1+it) \neq 0$（$t \neq 0$）须用另一方法，即 Hadamard（1896）的三角不等式；其前提是 $\zeta(s)$ 已在 $\sigma = 1$ 附近由解析延拓定义，见第~\ref{ch:zero_distribution}~章。

%==============================================================
\section{\texorpdfstring{$\zeta(s)$}{Zeta(s)} 的实部与虚部及其解析性}
\label{sec:holo}
%==============================================================

为便于后续三维几何形态分析，我们将 Dirichlet 级数形式的 $\zeta(s)$ 分解为实部与虚部。几何形态分析将函数分拆为实曲面、虚曲面和模曲面进行，以讨论各自特点及其相互关系。

当 $\operatorname{Re}(s) = \sigma > 1$ 时，
\begin{equation}
\zeta(\sigma + it) =\sum_{n=1}^{\infty} n^{-(\sigma + it)} 
=\sum_{n=1}^{\infty} n^{-\sigma} e^{-it \log n}
=\sum_{n=1}^{\infty} \frac{\cos(t \log n)}{n^{\sigma}} 
- i\sum_{n=1}^{\infty} \frac{\sin(t \log n)}{n^{\sigma}}.
\end{equation}
因此，实部与虚部分别为
\begin{align}
\mathrm{Re}[\zeta(\sigma + it)] 
&=\sum_{n=1}^{\infty} \frac{\cos(t \log n)}{n^{\sigma}}, \\
\mathrm{Im}[\zeta(\sigma + it)] 
&= -\sum_{n=1}^{\infty} \frac{\sin(t \log n)}{n^{\sigma}}.
\end{align}
注意：当 $n=1$ 时，$\log 1 = 0$，故 $\sin(t \log 1) = 0$，该项对虚部无贡献。

有了实部、虚部的级数表示，便可直接验证 Cauchy--Riemann 方程，从而说明 $\zeta(s)$ 在 $\operatorname{Re}(s)>1$ 内全纯。

\begin{proof}
记 $u(\sigma, t) = \mathrm{Re}[\zeta(\sigma + it)]$，$v(\sigma, t) = \mathrm{Im}[\zeta(\sigma + it)]$。由于级数及其逐项求导后的级数在 $\sigma > 1$ 的任意有界闭子集（如 $\sigma \geq 1+\epsilon$）上绝对且一致收敛，我们可在求和号下逐项求导。

对实部求偏导：
\[
\frac{\partial u}{\partial \sigma} =\sum_{n=1}^{\infty} \frac{\partial}{\partial \sigma} \left( n^{-\sigma} \cos(t \log n) \right) 
= -\sum_{n=1}^{\infty} (\log n) \, n^{-\sigma} \cos(t \log n),
\]
\[
\frac{\partial u}{\partial t} =\sum_{n=1}^{\infty} \frac{\partial}{\partial t} \left( n^{-\sigma} \cos(t \log n) \right) 
= -\sum_{n=1}^{\infty} (\log n) \, n^{-\sigma} \sin(t \log n).
\]

对虚部求偏导：
\[
\frac{\partial v}{\partial \sigma} = -\sum_{n=1}^{\infty} \frac{\partial}{\partial \sigma} \left( n^{-\sigma} \sin(t \log n) \right) 
=\sum_{n=1}^{\infty} (\log n) \, n^{-\sigma} \sin(t \log n),
\]
\[
\frac{\partial v}{\partial t} = -\sum_{n=1}^{\infty} \frac{\partial}{\partial t} \left( n^{-\sigma} \sin(t \log n) \right) 
= -\sum_{n=1}^{\infty} (\log n) \, n^{-\sigma} \cos(t \log n).
\]

由此直接验证 Cauchy--Riemann 方程：
\[
\frac{\partial u}{\partial \sigma} = \frac{\partial v}{\partial t}, \qquad 
\frac{\partial u}{\partial t} = -\frac{\partial v}{\partial \sigma}.
\]
偏导数连续（因级数在 $\sigma > 1$ 的任意有界闭子集上一致收敛），故 $\zeta(s)$ 在 $\operatorname{Re}(s) > 1$ 内为全纯函数。
\end{proof}


%==============================================================
\section{Euler 乘积与 Dirichlet 级数的生成函数}
\label{subsec:arith_dirichlet_gen}
\label{subsec:chebyshev_zeta_relation}
\label{sec:euler_dirichlet_gen}
%==============================================================

\paragraph{为何需要生成函数}
第~\ref{ch:elementary_number_theoretic_functions}~章关心的是\textbf{部分和}：
\[
  M(x)=\sum_{n\le x}\mu(n),\qquad
  \psi(x)=\sum_{n\le x}\Lambda(n),\qquad
  J(x)=\sum_{p^k\le x}\frac1k
  =\sum_{n\le x}\frac{\Lambda(n)}{\log n}.
\]
它们是阶梯函数，定义清楚，但暂时只是「算术求和」。

第~\ref{ch:integral_transforms}~章的 Perron 积分公式则说：若已知系数 $\{a_n\}$ 的
Dirichlet 级数
\[
  F(s)=\sum_{n=1}^{\infty}a_n n^{-s}
  \qquad(\operatorname{Re}s\text{ 充分大}),
\]
则部分和可写成竖线积分
\[
  A(x)=\sum_{n\le x}a_n
  =\frac{1}{2\pi i}\int_{(c)} F(s)\frac{x^s}{s}\,\mathrm{d}s.
\]
因此，要把 $M$、$\psi$、$J$ 送进复分析，\textbf{关键一步}是弄清：
当 $a_n=\mu(n)$、$\Lambda(n)$ 或 $\Lambda(n)/\log n$ 时，$F(s)$ 等于什么。

本节要做的，正是借助本章已有的 \textbf{Euler 乘积}，在 $\operatorname{Re}s>1$ 上证明
\begin{align*}
  \sum\mu(n)n^{-s}&=\frac1{\zeta(s)},&
  \sum\Lambda(n)n^{-s}&=-\frac{\zeta'(s)}{\zeta(s)},&
  \sum\frac{\Lambda(n)}{\log n}\,n^{-s}&=\log\zeta(s).
\end{align*}
这些等式称为相应系数的 \textbf{Dirichlet 生成函数}：
左边是算术系数的级数，右边是 $\zeta$ 的解析表达式。
有了它们，$F$ 就具体了；\textbf{本节只证明这些等式}，
不把 $M,\psi,J$ 真正写成积分（那只需把已证的 $F$ 代入上面的 Perron 积分公式，
见第~\ref{ch:discrete_continuous}~章）。

证明只用 Euler 乘积与无零点，不再重证第二章对 $\mu$、$\Lambda$ 的定义与反演。

\subsection{三条对应关系}

在 $\operatorname{Re}s>1$ 上，目标恒等式为：

\begin{center}
\renewcommand{\arraystretch}{1.55}
{\small
\begin{tabular}{|c|c|c|c|}
\hline
\textbf{系数 $a_n$} &
\textbf{部分和 $A(x)$} &
\textbf{Dirichlet 级数 $F(s)$} &
\textbf{等于} \\
\hline
$\mu(n)$ &
$M(x)$ &
$\displaystyle\sum\mu(n)n^{-s}$ &
$1/\zeta(s)$ \\
\hline
$\Lambda(n)$ &
$\psi(x)$ &
$\displaystyle\sum\Lambda(n)n^{-s}$ &
$-\zeta'(s)/\zeta(s)$ \\
\hline
$\Lambda(n)/\log n\ (n\ge 2)$ &
$J(x)$ &
$\displaystyle\sum\dfrac{\Lambda(n)}{\log n}\,n^{-s}$ &
$\log\zeta(s)$ \\
\hline
\end{tabular}}
\end{center}

\vspace{0.4em}
\noindent
表中的系数 $\mu(n)$、$\Lambda(n)$ 及 $\Lambda(n)/\log n$
一律沿用第~\ref{ch:elementary_number_theoretic_functions}~章的定义；
下面给出最简回顾（细节与数值表以该章为准）：
\begin{itemize}
  \item $\mu(1)=1$；若 $n$ 为 $k$ 个不同素数之积则 $\mu(n)=(-1)^k$；
        含平方因子则 $\mu(n)=0$（第~\ref{sec:mobius_function}~节）。
  \item $\Lambda(n)=\log p$ 当 $n=p^k$（$k\ge 1$），否则 $0$
        （第~\ref{sec:von_mangoldt}~节）；
        从而在素数幂处 $\Lambda(n)/\log n=1/k$。
  \item 二者另有初等卷积 $\Lambda(n)=-\sum_{d\mid n}\mu(d)\log d$
        （第~\ref{subsec:Lambda_mu_conv}~节），与 $\zeta$ 无关；
        本节\textbf{不}用该卷积证明生成函数，只用 Euler 乘积。
\end{itemize}

$\sum\mu n^{-s}=1/\zeta$ 与 $\sum\Lambda n^{-s}=-\zeta'/\zeta$
\textbf{彼此独立}：前者是 Euler 乘积取倒数再展开，后者是对 Euler 乘积取对数再求导；
证明 $\Lambda$ 的生成函数不需要 $\mu$，反之亦然。
$\log\zeta$ 则是取对数、尚未对 $s$ 求导的中间式，系数恰为 $\Lambda/\log n$。

\subsection{\texorpdfstring{$\mu$}{mu} 的生成函数}

\begin{theorem}{Möbius 函数的 Dirichlet 生成函数}{mu_dirichlet}
对 $\operatorname{Re}(s) > 1$，有
\begin{equation}
\sum_{n=1}^{\infty} \frac{\mu(n)}{n^{s}} = \frac{1}{\zeta(s)}.
\label{eq:mu_dirichlet}
\end{equation}
其中系数 $\mu(n)$ 取自第~\ref{ch:elementary_number_theoretic_functions}~章第~\ref{sec:mobius_function}~节。
\end{theorem}

\begin{proof}
由 Euler 乘积（定理~\ref{thm:euler_product}），当 $\operatorname{Re}(s)>1$ 时
\[
\zeta(s) = \prod_{p} \frac{1}{1-p^{-s}},
\]
故
\[
\frac{1}{\zeta(s)} = \prod_{p} (1-p^{-s}).
\]
将右侧乘积展开：每个无平方因子的 $n=\prod p_i$ 贡献 $(-1)^{\omega(n)}$，
含平方因子的 $n$ 贡献 $0$。这与第~\ref{sec:mobius_function}~节对 $\mu(n)$ 的定义逐项一致，
故
\[
\prod_{p} (1-p^{-s}) = \sum_{n=1}^{\infty} \frac{\mu(n)}{n^{s}}.
\]
绝对收敛允许重排。
\end{proof}

该式在 $\operatorname{Re}s>1$ 上是 Euler 乘积的直接代数后果。
由无零点定理，该半平面内 $\zeta\neq 0$，故 $1/\zeta$ 全纯。
$M(x)=\sum_{n\le x}\mu(n)$ 正是左侧级数的部分和；
第~\ref{ch:discrete_continuous}~章将 $1/\zeta$ 代入 Perron 积分公式，即回到 $M(x)$。

延拓之后可读奇性对应（预告）：若 $\zeta(\rho)=0$（计重数），则 $1/\zeta$ 在 $\rho$ 有极点；
若 $\zeta$ 在 $s=1$ 有单极点，则 $1/\zeta$ 在 $s=1$ 有零点——这些点均不在 $\operatorname{Re}s>1$ 内。

\subsection{\texorpdfstring{$\Lambda$}{Lambda} 的生成函数与 \texorpdfstring{$\log\zeta$}{log zeta}}

\begin{theorem}{von Mangoldt 函数的 Dirichlet 生成函数}{Lambda_dirichlet}
对 $\operatorname{Re}(s) > 1$，有
\begin{equation}
\sum_{n=1}^{\infty} \frac{\Lambda(n)}{n^{s}} = -\frac{\zeta'(s)}{\zeta(s)}.
\label{eq:Lambda_dirichlet}
\end{equation}
其中系数 $\Lambda(n)$ 取自第~\ref{ch:elementary_number_theoretic_functions}~章第~\ref{sec:von_mangoldt}~节。
\end{theorem}

\begin{proof}
对 Euler 乘积取对数：
\[
\log\zeta(s) = -\sum_{p} \log(1-p^{-s}) = \sum_{p} \sum_{k=1}^{\infty} \frac{p^{-ks}}{k}.
\]
两边对 $s$ 求导得
\[
\frac{\zeta'(s)}{\zeta(s)}
= -\sum_{p} \sum_{k=1}^{\infty} \frac{\log p}{p^{ks}}
= -\sum_{n=1}^{\infty} \frac{\Lambda(n)}{n^{s}},
\]
其中最后一步引用第~\ref{sec:von_mangoldt}~节：$\Lambda(n)$ 仅在 $n=p^k$ 处取值 $\log p$，其余为 $0$。
（双重级数在 $\operatorname{Re}s>1$ 上绝对收敛，故可逐项求导。）
\end{proof}

取对数、尚未求导的中间式给出 $\log\zeta$ 的 Dirichlet 级数（$\operatorname{Re}s>1$）：
\begin{equation}
\log\zeta(s)
= \sum_{p}\sum_{k\ge 1}\frac{p^{-ks}}{k}
= \sum_{n=2}^{\infty}\frac{\Lambda(n)}{\log n}\,n^{-s},
\label{eq:log_zeta_dirichlet}
\end{equation}
右侧系数在 $n=p^k$ 处为 $1/k=\Lambda(n)/\log n$，与 $J$ 的跳跃一致；
故 $J(x)$ 是该级数的部分和。

\subsection{三条恒等式的汇总}

合写为
\begin{equation}
\frac{1}{\zeta(s)} = \sum_{n=1}^{\infty} \frac{\mu(n)}{n^{s}}, \qquad
-\frac{\zeta'(s)}{\zeta(s)} = \sum_{n=1}^{\infty} \frac{\Lambda(n)}{n^{s}}, \qquad
\log\zeta(s) = \sum_{n=2}^{\infty}\frac{\Lambda(n)}{\log n}\,n^{-s}
\qquad (\operatorname{Re}s>1).
\label{eq:mu_Lambda_pair}
\end{equation}
第~\ref{ch:discrete_continuous}~章将 $F=1/\zeta$、$-\zeta'/\zeta$、$\log\zeta$
分别代入 Perron 积分公式，得到 $M$、$\psi$、$J$ 在 $\operatorname{Re}s>1$ 上的竖线积分表示。

对数导数 $\zeta'/\zeta$ 的极点位置可直接由定义读出（延拓后的语言，预告）：
在 $\zeta$ 的 $m$ 阶零点处，$\zeta'/\zeta$ 有一阶极点、留数 $m$；
在 $\zeta$ 的 $n$ 阶极点处（本书中即 $s=1$ 的单极点），$\zeta'/\zeta$ 有一阶极点、留数 $-n$。
第~\ref{ch:riemann_explicit_formula}~章对 $-\zeta'/\zeta$ 作围道积分时，留数项正来自这些点。

%==============================================================
\section{黎曼 \texorpdfstring{$\zeta$}{Zeta} 函数的对称性}
\label{sec:schwarz}
\label{sec:zeta_symmetry}
%==============================================================

本节在 $\operatorname{Re}s>1$ 上建立
\[
  \zeta(\bar{s})=\overline{\zeta(s)},
\]
所用工具是第~\ref{ch:complex_analysis_intro}~章第~\ref{sec:schwarz_reflection}~节的
Schwarz 反射原理；
并给出不依赖该原理、仅用级数逐项共轭的直接验证。

\subsection{实轴上取实值}

当 $s=\sigma$ 为大于 $1$ 的实数时，Dirichlet 级数的每一项 $n^{-\sigma}$ 均为正实数，故
\begin{equation}
  \zeta(\sigma)
  =\sum_{n=1}^{\infty}\frac{1}{n^{\sigma}}
  \in\mathbb{R},
  \qquad \sigma>1.
  \label{eq:zeta_real_on_real}
\end{equation}
这是 Schwarz 反射的前提条件：函数在实轴段上连续且取实值。

\subsection{由 Schwarz 反射原理推出共轭对称}

\begin{theorem}{黎曼 $\zeta$ 函数的共轭对称性（$\operatorname{Re}s>1$）}{zeta_symmetry}
\label{cor:zeta_symmetry}
对一切 $\operatorname{Re}s>1$，有
\begin{equation}
  \zeta(\bar{s})
  =\overline{\zeta(s)}.
  \label{eq:zeta_symmetry}
\end{equation}
\end{theorem}

\begin{proof}
取
\[
  D^{+}
  =\bigl\{s:\operatorname{Im}s>0,\;\operatorname{Re}s>1\bigr\},
  \qquad
  I=(1,+\infty),
  \qquad
  D^{-}
  =\bigl\{s:\operatorname{Im}s<0,\;\operatorname{Re}s>1\bigr\}.
\]
由本章前面的全纯性结论，$\zeta$ 在 $D^{+}$ 内全纯；
由~\eqref{eq:zeta_real_on_real}，它在 $I$ 上连续且取实值。
第~\ref{ch:complex_analysis_intro}~章第~\ref{sec:schwarz_reflection}~节（Schwarz 反射原理）给出：
$\zeta$ 可关于实轴对称延拓至 $D^{+}\cup I\cup D^{-}$，且延拓满足
\[
  \zeta(\bar{s})=\overline{\zeta(s)}
  \qquad\bigl(s\in D^{+}\cup I\cup D^{-}\bigr).
\]
而 $D^{+}\cup I\cup D^{-}$ 恰为半平面 $\operatorname{Re}s>1$，故得~\eqref{eq:zeta_symmetry}。
\end{proof}

\paragraph{级数直接验证（不经过反射原理）}
因系数 $a_n=1$ 为实数，且当 $\operatorname{Re}s=\sigma>1$ 时级数绝对收敛，可对级数逐项取复共轭：
\[
  \overline{\zeta(\sigma+it)}
  =\overline{\sum_{n=1}^{\infty}n^{-(\sigma+it)}}
  =\sum_{n=1}^{\infty}n^{-(\sigma-it)}
  =\zeta(\sigma-it)
  =\zeta(\overline{\sigma+it}).
\]
这与~\eqref{eq:zeta_symmetry} 一致，并说明在绝对收敛半平面内，
共轭对称也可不引用 Schwarz 反射而直接读出。
Schwarz 路线的价值在于：它把「实轴取实值 $\Rightarrow$ 共轭对称」写成一般原理；
后文经解析延拓后的 $\zeta$ 在更大区域上仍满足同一恒等式时，仍可沿同一思路论证
（延拓后的全平面情形见第~\ref{ch:riemann_zeta_continuation}~章之后）。

\subsection{对实部、虚部与模的推论}

由~\eqref{eq:zeta_symmetry}，在 $\operatorname{Re}s=\sigma>1$ 上立刻得到：
\begin{enumerate}[label=(\roman*)]
  \item 实部 $\operatorname{Re}\zeta(\sigma+it)$ 是 $t$ 的\textbf{偶}函数，
        虚部 $\operatorname{Im}\zeta(\sigma+it)$ 是 $t$ 的\textbf{奇}函数；
  \item 模 $|\zeta(\sigma+it)|$ 关于 $t=0$ 偶对称；
  \item 因而后文三维曲面中，实部曲面与模曲面关于平面 $t=0$ 对称，
        虚部曲面关于该平面反对称。
\end{enumerate}
几何含义：共轭点 $s$ 与 $\bar{s}$ 映为共轭点，像曲线关于 $w$ 平面实轴对称
（几何示意见第~\ref{ch:visualization}~章）。


%==============================================================
\section{Dirichlet 级数形式的几何示意（截断示意）}
\label{sec:zeta_visualization}
%==============================================================

本节用\textbf{截断级数}
\begin{equation}
  \zeta_N(s)=\sum_{n=1}^{N}n^{-s}
  \qquad(N=200)
\end{equation}
绘制实部、虚部与模的三维曲面，目的是直观看见 $\operatorname{Re}s>1$ 上级数定义的 $\zeta$
在逼近直线 $\sigma=1$ 时的形态，并分清\textbf{真性质}与\textbf{截断伪像}。
图为示意，\textbf{不代替}无零点、延拓等证明。

\subsection{截断误差与梳状结构}
\label{sec:trunc-error}

当 $\sigma>1$ 时，
\[
  \bigl|\zeta(s)-\zeta_N(s)\bigr|
  \le\sum_{n=N+1}^{\infty}n^{-\sigma}
  \le\frac{N^{1-\sigma}}{\sigma-1}.
\]
$\sigma$ 固定、$N$ 大时误差 $O(N^{1-\sigma})\to 0$；
$\sigma\to 1^+$ 时上界变差，截断影响被放大。

\paragraph{梳状结构来自高阶项截断}
绘图中 $\sigma\to 1^+$ 一侧常见的\textbf{梳状峰谷}，主要不是「无穷级数 $\zeta$ 的固有周期结构」，
而是有限和 $\zeta_N$ 的干涉图案：
\begin{itemize}
  \item 实部 $\sum_{n=1}^{N}\cos(t\log n)\,n^{-\sigma}$、虚部 $-\sum\sin(t\log n)\,n^{-\sigma}$
        是频率集 $\{\log n\}_{n\le N}$ 上的有限叠加；
  \item 最大项 $n=N$ 的特征频率为 $\log N$，梳齿间距约为
        \[
          \Delta t=\frac{2\pi}{\log N}
          \approx\frac{2\pi}{\log 200}\approx 1.19
          \qquad(N=200);
        \]
  \item $\sigma$ 接近 $1$ 时 $n^{-\sigma}$ 衰减变慢，\textbf{高阶项}（大 $n$）相对增强，
        截断后的相长/相消干涉更醒目；$\sigma$ 较大时高阶项被压住，曲面趋平。
\end{itemize}
因而：梳状结构随 $N$ 改变而变（间距 $\propto 1/\log N$），属于\textbf{级数高阶项被截断}带来的数值效应；
$\sigma\le 1$ 处更不能把 $\zeta_N$ 的振荡读成延拓后 $\zeta$ 的解析结构。
$N\to\infty$ 且 $\sigma>1$ 固定时，$\zeta_N\to\zeta$，截断梳齿在理论极限下消隐于真函数。

\subsection{三曲面要点（观察）}

高度经 $\arctan(2\,\cdot\,)$ 压缩。实部、虚部、模见图~\ref{fig:real}--\ref{fig:Modular}。

\begin{itemize}
  \item \textbf{共同}：$\sigma$ 大时平坦；$\sigma\to 1^+$ 时抬升，并出现上述截断梳状起伏。
  \item \textbf{实部}：关于 $t$ 偶对称；绘图域内多在 $z=0$ 上方（观察）。
        真 $\zeta$ 的 $\operatorname{Re}s>1$ 无零点由 Euler 乘积保证，图不能代替证明。
  \item \textbf{虚部}：关于 $t$ 奇对称（$\sin$ 的奇性 / 共轭对称）；
        实轴上虚部为 $0$；尖峰可正可负。
  \item \textbf{模}：非负、关于 $t$ 偶对称；绘图域内不触底，与 $\sigma>1$ 无零点方向一致（观察）。
        $t=0$、$\sigma\to 1^+$ 的高脊对应 $\zeta(\sigma)\to+\infty$（$s=1$ 处级数发散 / 极点预兆）。
\end{itemize}

\begin{figure}[htbp]
  \centering
  \includegraphics[width=0.75\textwidth]{全书图形/5-Z实曲面.pdf}
  \caption{截断 Dirichlet 级数 $\zeta_N$ 的实部曲面（示意）。$\sigma\to 1^+$ 邻域的梳状起伏主要来自有限 $N$ 的高阶项截断；与 $\sigma>1$ 无零点一致处仅为观察。}
  \label{fig:real}
\end{figure}

\begin{figure}[htbp]
  \centering
  \includegraphics[width=1.0\textwidth]{全书图形/5-Z虚曲面.pdf}
  \caption{截断级数 $\zeta_N$ 的虚部曲面：关于 $t=0$ 奇对称；左侧密集峰谷同属截断干涉。}
  \label{fig:Imaginary}
\end{figure}

\begin{figure}[htbp]
  \centering
  \includegraphics[width=1.0\textwidth]{全书图形/5-Z模曲面.pdf}
  \caption{截断级数 $\zeta_N$ 的模曲面：非负；绘图域内无触底（观察）；$s\to 1^+$ 脊线为发散预兆。}
  \label{fig:Modular}
\end{figure}

\begin{center}
\renewcommand{\arraystretch}{1.35}
{\small
\begin{tabular}{|l|l|l|l|}
\hline
& 实部 & 虚部 & 模 \\
\hline
$t$ 对称 & 偶 & 奇 & 偶 \\
$\sigma\to 1^+$ & 抬升 $+$ 截断梳状 & 双向截断尖峰 & 高脊 $+$ 截断起伏 \\
与理论 & 无零点：Euler 乘积 & 奇对称：共轭/$\sin$ & 无触底：观察$\parallel$无零点 \\
\hline
\end{tabular}
}
\end{center}

\vspace{0.5em}
\noindent
第~\ref{ch:riemann_zeta_continuation}~章将在相同精神下给出延拓后的对比曲面。
边界 $\zeta(1+it)\neq 0$ 的证明见第~\ref{ch:zero_distribution}~章（Hadamard 等），
不能由图上「模不触底」代替。

%==============================================================
\section{Dirichlet 级数的局限性：解析延拓的必要性}
\label{sec:limitations}
%==============================================================

前几节已建立 $\zeta(s)$ 在 $\operatorname{Re}(s) > 1$ 内的定义与基本性质
（含 Euler 乘积、无零点、全纯、对称、$\mu$/$\Lambda$ 生成函数，以及截断示意）。
本节从\textbf{以 Dirichlet 级数作为 $\zeta$ 的原初定义}这一全局立场，
归纳该框架覆盖不了什么，从而说明解析延拓的必要性。

与第~\ref{sec:euler-fail}~节的分工：该节只讨论\textbf{乘积工具}在 $\sigma=1$ 上的方法边界；
本节讨论\textbf{级数定义本身}的定义域与目标对象之间的缺口。

\medskip
\textbf{局限一：$\operatorname{Re}(s)\le 1$ 上级数不能定义 $\zeta(s)$。}

章首已得 $\sigma_a=\sigma_c=1$：$\operatorname{Re}s>1$ 时级数绝对收敛并可定义 $\zeta(s)$；
$s=1$ 时化为调和级数而发散。因而 $\operatorname{Re}s\le 1$ 上不能用本级数定义 $\zeta$。
研究素数分布所需的非平凡零点（$0<\operatorname{Re}s<1$）
以及特殊值 $\zeta(0)=-\frac12$、$\zeta(-1)=-\frac1{12}$ 等
均落在级数定义域之外，Dirichlet 级数无法覆盖。

\medskip
\textbf{局限二：直线 $\operatorname{Re}s=1$ 上的无零点证不出，而素数定理需要它。}

分清三层事实：

\begin{enumerate}[label=(\roman*)]
  \item \textbf{本章已证}：当 $\operatorname{Re}s>1$ 时 $\zeta(s)\neq 0$。
        证明用 Euler 乘积：每个因子 $(1-p^{-s})^{-1}\neq 0$，且乘积绝对收敛故整体非零。
  \item \textbf{素数定理需要更强的结论}：对一切实数 $t\neq 0$，
        \[
          \zeta(1+it)\neq 0.
        \]
        即 $\zeta$ 在收敛边界直线 $\operatorname{Re}s=1$ 上（除可能的 $s=1$ 外）也无零点。
        没有它，经典的解析证明推不出 $\psi(x)\sim x$（见第~\ref{ch:prime_number_theorem}~章）。
  \item \textbf{本章工具推不到 (ii)}：
        \begin{itemize}
          \item 点 $1+it$（$t\neq 0$）满足 $\operatorname{Re}s=1$，\textbf{不在}级数的绝对收敛半平面内；
                本章用级数定义的 $\zeta$ 根本\textbf{没有}在这些点上取值，
                谈不上「$\zeta(1+it)$ 是否为零」。
          \item Euler 乘积在 $\sigma=1$ 上\textbf{不再绝对收敛}（第~\ref{sec:euler-fail}~节，Mertens $\sum 1/p=\infty$），
                不能把「因子非零 $\Rightarrow$ 乘积非零」的论证原样搬到边界上。
        \end{itemize}
\end{enumerate}

因此：从「$\operatorname{Re}s>1$ 无零点」到「$\zeta(1+it)\neq 0$」不是同一定理的小推广，
而是\textbf{定义域与方法都要升级}——须先将 $\zeta$ 解析延拓到 $\sigma=1$ 附近，
再用 Hadamard 三角不等式等另行证明
（第~\ref{ch:zero_distribution}~、\ref{ch:prime_number_theorem}~章）。

\medskip
\textbf{局限三：函数方程在级数框架内无从建立。}

第~\ref{ch:riemann_zeta_continuation}~章将证明黎曼函数方程
（形状预告，本章不使用其成立性）
\[
  \zeta(s) = 2^s \pi^{s-1} \sin\!\left(\frac{\pi s}{2}\right)
             \Gamma(1-s)\, \zeta(1-s).
\]
它连接 $\sigma>1$ 与 $\sigma<0$。
当 $\operatorname{Re}s>1$ 时 $\operatorname{Re}(1-s)<0$，$\zeta(1-s)$ 不能再用 Dirichlet 级数定义，
故该式无法在本章的级数框架内建立，而依赖于全平面的亚纯延拓。

\medskip
\textbf{局限四：截断级数的数值伪像。}

由第~\ref{sec:trunc-error}~节：有限和 $\zeta_N$ 在 $\sigma\to 1^+$ 附近的\textbf{梳状结构}
主要来自\textbf{高阶项截断}的干涉，并非无穷级数 $\zeta$ 的固有周期几何；
$\sigma\le 1$ 处更不能把 $\zeta_N$ 读成延拓后的 $\zeta$。
图上的「无触底」「无交线」只是观察，不能代替证明。

\medskip
综上：Dirichlet 级数在 $\sigma>1$ 上给出了全纯的 $\zeta(s)$ 与 Euler 乘积，
但定义域、边界无零点、函数方程与特殊值都要求走出该级数框架——
这正是后文解析延拓的出发点。

%==============================================================
\section{积分表示与 Dirichlet 级数的等价性}\label{sec:integral-representation-eq}
%==============================================================

尽管 Dirichlet 级数有其固有的局限，但通过积分表示（Mellin 变换形式），我们可以在 $\operatorname{Re}(s)>1$ 区域内建立一个等价形式。

建立该等价性需要 \textbf{$\Gamma$ 函数} $\Gamma(s)$。其复解析性质、因式分解与全平面延拓见第~\ref{ch:gamma_function}~章；本节为推导 $\zeta(s)$ 的积分表示，先给出 $\Gamma(s)$ 在 $\operatorname{Re}(s)>0$ 上的经典积分定义：

\begin{definition}{$\Gamma$ 函数在右半平面上的积分表示}{}
对于满足 $\operatorname{Re}(s)>0$ 的复数 $s$，定义如下收敛的无穷反常积分：
\begin{equation}
    \Gamma(s) = \int_0^\infty x^{s-1}e^{-x}\,\mathrm{d}x, \qquad \operatorname{Re}(s)>0.
    \label{eq:gamma-def-ch4}
\end{equation}
通过分部积分，当 $s=n$ 为正整数时，该积分退化为阶乘形式：$\Gamma(n)=(n-1)!$。
\end{definition}

由此可写出 $\zeta(s)$ 在 $\operatorname{Re}(s)>1$ 内的经典积分表示。该式本身仍限于 $\operatorname{Re}(s)>1$，尚未延拓；它是后续路径变形的出发点：
\begin{equation}
    \zeta(s) = \frac{1}{\Gamma(s)} \int_0^\infty \frac{x^{s-1}}{e^x-1}\,\mathrm{d}x,
    \qquad \operatorname{Re}(s)>1.
    \label{eq:zeta-integral}
\end{equation}

这一表示与 Dirichlet 级数的等价性可以通过初等运算直接验证。将被积函数中的 \(\frac{1}{e^x-1}\) 展开为几何级数：
\[
\frac{1}{e^x-1} = \frac{e^{-x}}{1-e^{-x}} = \sum_{n=1}^\infty e^{-nx}, \qquad x>0.
\]
代入 \eqref{eq:zeta-integral}，并在绝对收敛的条件下交换求和与积分：
\[
\int_0^\infty \frac{x^{s-1}}{e^x-1}\,\mathrm{d}x
= \sum_{n=1}^\infty \int_0^\infty x^{s-1} e^{-nx}\,\mathrm{d}x.
\]
对每个积分作变量代换 \(u = nx\)，得
\[
\int_0^\infty x^{s-1} e^{-nx}\,\mathrm{d}x = \frac{1}{n^s} \int_0^\infty u^{s-1} e^{-u}\,\mathrm{d}u = \frac{\Gamma(s)}{n^s}.
\]
因此
\[
\int_0^\infty \frac{x^{s-1}}{e^x-1}\,\mathrm{d}x = \Gamma(s) \sum_{n=1}^\infty \frac{1}{n^s} = \Gamma(s)\,\zeta(s),
\]
两边除以 \(\Gamma(s)\) 即得 \eqref{eq:zeta-integral}。由此可见，在 \(\operatorname{Re}(s)>1\) 内，积分表示与 Dirichlet 级数表示完全等价。这一等价性为后续通过积分路径变形将 \(\zeta(s)\) 解析延拓至全平面提供了出发点，但具体延拓方法将在第~\ref{ch:riemann_zeta_continuation}~章（$\zeta$ 函数的解析延拓与函数方程）中详述。

\begin{center}
\fbox{\parbox{0.92\textwidth}{\centering
\vspace{0.2em}
\textbf{第五章：重要定理与核心公式汇总}\\[6pt]
\renewcommand{\arraystretch}{1.6}
\begin{tabular}{lp{9.5cm}}
\hline
\textbf{名称} & \textbf{数学内容 / 公式表达} \\
\hline
$\zeta(s)$ 的 Dirichlet 级数定义 & $\zeta(s)=\displaystyle\sum_{n=1}^{\infty}\dfrac{1}{n^s}$，$\Re(s)>1$ \\
Euler 乘积公式 & $\zeta(s)=\displaystyle\prod_p(1-p^{-s})^{-1}$，$\Re(s)>1$ \\
$\mu$ 的生成函数 & $\displaystyle\sum\mu(n)n^{-s}=1/\zeta(s)$，$\Re(s)>1$ \\
$\Lambda$ 的生成函数 & $\displaystyle\sum\Lambda(n)n^{-s}=-\zeta'/\zeta(s)$，$\Re(s)>1$ \\
收敛域与全纯性 & Dirichlet 级数与 Euler 乘积在 $\Re(s)>1$ 绝对收敛，$\zeta(s)$ 在该半平面全纯 \\
非零性 & $\zeta(s)\neq 0$，$\Re(s)>1$（由 Euler 乘积保证） \\
共轭对称 & $\zeta(\bar{s})=\overline{\zeta(s)}$，$\Re(s)>1$（Schwarz 反射 / 逐项共轭） \\
极点预兆 & $\sigma\to1^+$ 时级数发散，预示 $s=1$ 处一阶极点 \\
\hline
\end{tabular}
\vspace{0.2em}
}}
\end{center}

\section*{本章小结与后续展望}
\addcontentsline{toc}{section}{本章小结与后续展望}

本章讨论了 $\zeta(s)$ 在绝对收敛域 $\operatorname{Re}(s)>1$ 内的定义与解析性质。
主线为 Dirichlet 级数、Euler 乘积、无零点与全纯性；
在此之后由 Euler 乘积建立 $\mu$、$\Lambda$ 的生成函数
$1/\zeta$、$-\zeta'/\zeta$ 与 $\log\zeta$。
共轭对称与截断级数示意进一步说明 $\operatorname{Re}s>1$ 上的结构与逼近 $\sigma=1$ 时的边界。

下一章（第~\ref{ch:discrete_continuous}~章）将引用这些生成函数，
用第~\ref{ch:integral_transforms}~章的 Perron 积分公式把 $\psi,M,J$
（及经反演的 $\theta,\pi$）写成 $\operatorname{Re}s>1$ 上的竖线积分；
本章不讨论数论阶梯的积分表示。
再往后建立 $\Gamma$ 与解析延拓的一般理论，并在第~\ref{ch:riemann_zeta_continuation}~章
通过复围道积分和 Jacobi $\vartheta$ 函数模反演，将 $\zeta(s)$ 延拓至整个复平面。